% ===================================================================
%  TABELO N-RO 5 — La domo kaj ĝia konstruado.
%  Traduction 2026 d'apres texto/io, controlee sur texto/fr et
%  texto/es. Consigne : voir texto/eo/10-tabelo-01.tex.
%
%  Le tableau 5 est un DIALOGUE, et les deux volumes composent le nom
%  du parleur en petites capitales : on garde \textsc{}, que le rendu
%  laisse tomber en gardant le texte.
%
%  Les renvois du plan sont des LETTRES — (a) le couloir, (b) le
%  salon... — et non des numeros. On les ecrit comme l'ido les ecrit,
%  y compris la ou l'imprimeur a compose « (1) » pour « (l) ».
%
%  LES DEUX NOMS PROPRES SE TRADUISENT, ET L'IDO L'A DEJA FAIT :
%  Leblanc est « Blanko », Roux est « Rufo ». L'esperanto forme les
%  memes mots a partir des memes adjectifs — blanka, rufa — et n'a
%  donc rien a changer. C'est la seule colonne ou la question ne se
%  pose pas.
% ===================================================================

%%K t05-apar-1 apar
\VUsaut{6.0mm}
\VUtitre{15.0pt}{60}{TABELO N\textsuperscript{o} 5}
\VUsaut{1.4mm}
\VUtitre{11.4pt}{0}{La domo kaj ĝia konstruado.}
\VUsaut{2.2mm}

%%K t05-tit-1 sub
\VUpk{10.2pt}{40}{Bona renkonto.}
\VUsaut{3.0mm}

%%K t05-01-1 p
\textsc{Johano}. --- Onklo mia, mi tre deziras scii kiel oni konstruas
\VUgras{domon}.

%%K t05-01-2 p
\textsc{La onklo} (80). --- Tio estas tre facila, mia amiko, venu kun mi, mi
montros al vi tiun kiun sinjoro Bernardo (79) igas konstrui kaj kiun mi loĝos.

%%K t05-01-3 p
\textsc{Johano}. --- Kiu do desegnis la \VUgras{planon} \textsuperscript{(76)}
de ĉi tiu domo?

%%K t05-01-4 p
\textsc{La onklo}. --- \VUgras{La arkitekto} \textsuperscript{(75)}; li
prezentis tre klaran desegnaĵon kiu estis aprobita, kaj \VUgras{la
entreprenisto} \textsuperscript{(77)} komencis la laborojn.

%%K t05-01-5 p
\textsc{Johano}. --- Kiu do estas la entreprenisto?

%%K t05-01-6 p
\textsc{La onklo}. --- Li estas sinjoro Blanko, la patro de via amiko Jakobo.
Li direktas la laboristojn kun sinjoro Rufo, la arkitekto.

%%K t05-01-7 p
Nu! Jen ĝustatempe venas sinjoro Blanko kun sia \VUgras{mezurilo}
\textsuperscript{(78)}, kaj sinjoro Rufo kun sia plano.

%%K t05-01-8 p
Bonan tagon, sinjoroj. Kiel vi fartas?

%%K t05-01-9 p
\textsc{Sinjoro Rufo}. --- Tre bone, dankon. Bonan tagon, Johano, kiel kreskis
ĉi tiu infano!

%%K t05-01-10 p
\textsc{La onklo}. --- Jes, vere, li estas malgranda viro. Li estas tre
serioza, oni devas informi lin pri ĉio kion li vidas. Hodiaŭ, li volus scii
kiel oni konstruas domon.

%%K t05-tit-2 sub
\VUpk{10.2pt}{40}{La laboristoj.}
\VUsaut{3.0mm}

%%K t05-01-11 p
\textsc{Sinjoro Blanko}. --- Nu, venu kun mi, amiketo, mi tuj klarigos ĝin al
vi, ĉar mi tre amas la knabojn kiuj volas instrui sin.

%%K t05-01-12 p
Vi vidas tiun laboriston kiu tenas \VUgras{ŝovelilon} \textsuperscript{(82)} en
la mano: li estas \VUgras{terlaboristo} \textsuperscript{(81)}. Li elfosas
\VUgras{tranĉeon} \textsuperscript{(46)} por meti la akvotubaron. Li ankaŭ
elfosis la teron en la loko kie estas la domo, por ke la masonisto starigu la
kvar \VUgras{murojn}. La parto de tiuj muroj kiu estas en la tero estas
nomata \VUgras{fundamentoj} \textsuperscript{(2)}. La spaco enfermita inter la
fundamentoj formas la \VUgras{kelon} \textsuperscript{(3)}. Tiu kelo havas
malgrandan fenestron nomatan \VUgras{kelfenestro} \textsuperscript{(5)},
\VUgras{pordon} \textsuperscript{(4)} kaj ŝtuparon.

%%K t05-01-13 p
La \VUgras{masonisto} \textsuperscript{(108)} daŭrigis konstrui la murojn
lasante aperturojn por la \VUgras{pordoj} \textsuperscript{(8)} kaj la
\VUgras{fenestroj} \textsuperscript{(17)}.

%%K t05-01-14 p
\textsc{Johano}. --- Sed, tiun masoniston mi ne vidas!

%%K t05-01-15 p
S\textsuperscript{o} \textsc{Blanko}. --- Nun li finis labori en la domo. Li
estas apud la \VUgras{ĉevalejo} \textsuperscript{(54)}, sur sia
\VUgras{skafaldo} \textsuperscript{(110)}, li konstruas la muron de la
\VUgras{lavejo} \textsuperscript{(56)}.

%%K t05-01-16 p
Li havas trulon, trogon kaj \VUgras{plumbolinion} \textsuperscript{(109)}. Sed
vi ne memoros tiujn nomojn.

%%K t05-01-17 p
\textsc{Johano}. --- Certe jes, ĉar mi tuj petos de mia onklo malgrandan
trulon kaj trogon, kaj mi mem fabrikos plumbolinion per ŝnureto.

%%K t05-tit-3 sub
\VUsaut{3.0mm}
\VUpk{10.2pt}{40}{La materialo.}
\VUsaut{3.0mm}

%%K t05-01-18 p
\textsc{La onklo}. --- Ha! tre bone. Sed diru al mi, Johano, kion oni bezonas
por konstrui domon?

%%K t05-01-19 p
\textsc{Johano}. --- Oni bezonas \VUgras{kvadrojn} \textsuperscript{(59)} kaj
\VUgras{mortero} \textsuperscript{(65)}.

%%K t05-01-20 p
\textsc{La onklo}. --- Ĝuste, vidu tiun viron, kun granda ĉapelo, sur sia
\VUgras{ĉaro} \textsuperscript{(136)}. Li estas \VUgras{ĉarkondukisto}
\textsuperscript{(135)}. Ĉiutage li alkondukas ĉi tien \VUgras{ŝtonajn blokojn}
\textsuperscript{(60)}. Li malŝarĝas sian veturilon en la korto, kaj la
\VUgras{kvadrotranĉistoj} \textsuperscript{(83)} tranĉas la blokojn kaj segas
ilin per sia \VUgras{ŝtonsegilo} \textsuperscript{(85)}.
\VUgras{Manlaboristo} \textsuperscript{(146)} portas ilin al la masonisto kiu
lokas ilin.

%%K t05-01-21 p
Dume, la \VUgras{morterfaristo} \textsuperscript{(87)} preparas la morteron. Li
bezonas \VUgras{sablon} \textsuperscript{(62)}, \VUgras{kalkon}
\textsuperscript{(63)} kaj \VUgras{akvon} \textsuperscript{(64)}. La kalko
estas apud li en \VUgras{sako} \textsuperscript{(71)}, \VUgras{masonservisto}
\textsuperscript{(111)} alportas al li la sablon sur \VUgras{puŝĉareto}
\textsuperscript{(112)} kaj la \VUgras{metilernanto} \textsuperscript{(113)}
estas ĉe la pumpilo (58). Li plenigas \VUgras{sitelon} \textsuperscript{(114)}.
La mortero estas baldaŭ preta. La \VUgras{morterportisto}
\textsuperscript{(107)} prenas ĝin kaj supren portas, per la ŝtupetaro, sur la
skafaldon, kie laboras masonisto kiu finas tiun flankon de la domo.

%%K t05-01-22 p
Kaj nun, kiel estas nomata la laboristo kiu laboras la \VUgras{tegmenton}
\textsuperscript{(31)}?

%%K t05-01-23 p
\textsc{Johano}. --- Li estas \VUgras{tegmentisto} \textsuperscript{(118)}.

%%K t05-tit-4 sub
\VUsaut{3.0mm}
\VUpk{10.2pt}{40}{La tegmento de la domo.}
\VUsaut{3.0mm}

%%K t05-01-24 p
S\textsuperscript{o} \textsc{Blanko}. --- Jes, sed unue venas la
\VUgras{ĉarpentisto} \textsuperscript{(115)}.

%%K t05-01-25 p
La ĉarpentisto laboras la \VUgras{ĉarpentaĵon} \textsuperscript{(35)}. Li
kunigas la \VUgras{trabojn} \textsuperscript{(37)} per \VUgras{kejloj}
\textsuperscript{(117)}. Tiuj laboristoj, tie supre, kun siaj larĝaj velurbrakoj,
estas ĉarpentistoj. La unu tenas trabeton, kaj la alia segas kejlon per sia
\VUgras{hakilo} \textsuperscript{(116)}. Tiu kiu estas apud la
\VUgras{kameno} \textsuperscript{(41)} estas la tegmentisto. Li najlas latojn
sur la (*) traboj, kaj \VUgras{ardezojn} \textsuperscript{(66)} sur la latojn.
Kelkfoje li metas tegolojn.

%%K t05-noto-1 noto
\VUnotes{143.2mm}{%
(*) \VUgras{Balk-o} = G. \textit{Balken}, A. \textit{joist}, F. \textit{solive},
I. \textit{travicello}, R. \textit{balka}, H. Port. \textit{viga}. Teknika
radiko ĝis nun ne adoptita, sed proponinda por traduki la sencon de la vorto F.
\textit{solive}, kiu estas nek trabo F. \textit{poutre}, nek trabeto. Ĝi estas
ligna stango kvadratigita kiu subtenas plankon kaj plafonon.%
}

%%K t05-01-26 p
La tegmento de la ĉevalejo estas \VUgras{tegolkovrita} \textsuperscript{(55)},
tiu de la domo estas \VUgras{ardezkovrita} \textsuperscript{(32)} kaj tiu de la
verando estas \VUgras{zinka} \textsuperscript{(24)}.

%%K t05-01-27 p
\textsc{Johano}. --- Ĉu la tegmentisto laboras ankaŭ la zinkajn tegmentojn?

%%K t05-01-28 p
S\textsuperscript{o} \textsc{Blanko}. --- Ne, sed la \VUgras{zinklaboristo}
\textsuperscript{(121)} aŭ la \VUgras{plumblaboristo}, tiu kiu metas
\VUgras{lukon} \textsuperscript{(38)} sur la tegmenton de la domo. Li metas
ankaŭ la pluv-\VUgras{kanalojn} \textsuperscript{(43)} kaj la
\VUgras{tegmentkanalojn} \textsuperscript{(42)}. Li kunlutas la zinkon kaj la
ladon. Li fiksis la \VUgras{fulmoŝirmilon} \textsuperscript{(40)} kaj la
\VUgras{ventoflagon} \textsuperscript{(39)} sur la \VUgras{frontonon}
\textsuperscript{(36)}.

%%K t05-01-29 p
\textsc{Johano}. --- Tiel la domo estas finkonstruita?

%%K t05-tit-5 sub
\VUsaut{3.0mm}
\VUpk{10.2pt}{40}{La iloj.}
\VUsaut{3.0mm}

%%K t05-01-30 p
S\textsuperscript{o} \textsc{Blanko}. --- Ne komplete. La \VUgras{gipsisto}
\textsuperscript{(99)} konstruas per \VUgras{brikoj} \textsuperscript{(61)} la
vandojn kiuj dividas ĝin en diversajn ĉambrojn. Li ŝmiras per \VUgras{gipso}
\textsuperscript{(70)} la \VUgras{plafonojn} \textsuperscript{(26)} kaj la
murojn. Nun li laboras la plafonon de la \VUgras{verando}
\textsuperscript{(33)}. Li havas, kiel la masonisto, \VUgras{trulon}
\textsuperscript{(100)} kaj \VUgras{trogon} \textsuperscript{(101)}.

%%K t05-01-31 p
Poste la lignaĵisto laboras la pordojn, la fenestrojn, la \VUgras{parketojn}
\textsuperscript{(16)} kaj la \VUgras{ŝutrojn} \textsuperscript{(19)}.

%%K t05-01-32 p
Nun li laboras en la korto sur sia labor-\VUgras{benko} \textsuperscript{(90)}.
Li havas \VUgras{segilon} \textsuperscript{(94)} por segi la lignon (69),
\VUgras{rabotilon} \textsuperscript{(91)}, \VUgras{fiksilon}
\textsuperscript{(92)}, \VUgras{ĉizilon} \textsuperscript{(94 bis)},
\VUgras{cirkelon} \textsuperscript{(95)}, kaj lignan \VUgras{martelon}
\textsuperscript{(95 bis)}.

%%K t05-01-33 p
\textsc{Johano}. --- Ha! sed estas pli amuze lignaĵi ol masoni. Onklo mia, vi
aĉetos por mi laborbenkon, segilon kaj rabotilon, ĉar mi baldaŭ laboros
kabanon por Medoro.

%%K t05-01-34 p
\textsc{La onklo}. --- Jes, knabeto.

%%K t05-01-35 p
\textsc{Johano}. --- Kiel kontenta mi estas! Kaj kiu estas tiu kiu staras apud
la enirpordo?

%%K t05-tit-6 sub
\VUsaut{3.0mm}
\VUpk{10.2pt}{40}{La farbado.}
\VUsaut{3.0mm}

%%K t05-01-36 p
S\textsuperscript{o} \textsc{Blanko}. --- Li estas \VUgras{farbisto}
\textsuperscript{(102)}. Li staras sur sia ŝtupetaro kun poto da \VUgras{farbo}
\textsuperscript{(103)} kaj sia \VUgras{peniko} \textsuperscript{(104)}. Li
farbas la lignon de la enirpordo por pli longa konserviĝo.

%%K t05-01-37 p
Li ankaŭ gluas la paperojn en la apartamentoj.

%%K t05-01-38 p
La \VUgras{tapetisto} \textsuperscript{(107)} kiu estas sur \VUgras{ŝtupetaro}
\textsuperscript{(106)}, sur la \VUgras{balkono} \textsuperscript{(28)}, metas
\VUgras{kurtenon} \textsuperscript{(29)}.

%%K t05-01-39 p
La laboristo staranta apud la granda fenestro estas la \VUgras{vitristo}
\textsuperscript{(96)}. Li fiksas la vitro-\VUgras{platojn}
\textsuperscript{(18)} per \VUgras{mastiko} \textsuperscript{(73)}. Tiu alia
laboristo, genuiĝinta apud la pordo de la kelo estas \VUgras{seruristo}
\textsuperscript{(97)}. Li metas \VUgras{seruron} \textsuperscript{(6)}. Lia
\VUgras{fajlilo} \textsuperscript{(98)} estas apud li.

%%K t05-01-40 p
\textsc{Johano}. --- Ĉu estas ankaŭ seruristo, tiu laboristo kiu frapas per sia
\VUgras{martelo} \textsuperscript{(84)} sur feraĵon?

%%K t05-01-41 p
S\textsuperscript{o} \textsc{Blanko}. --- Pli ĝustadire, li estas forĝisto
(125). Li forĝas \VUgras{feraĵojn} \textsuperscript{(67)} por la
\VUgras{elektristo} \textsuperscript{(124)} kiu estas sur la tegmento de la
\VUgras{veturilejo} \textsuperscript{(51)}. Li havas \VUgras{forĝejon}
\textsuperscript{(126)}, \VUgras{amboson} \textsuperscript{(127)} kaj
\VUgras{prenilon} \textsuperscript{(128)}.

%%K t05-01-42 p
La elektristo fiksas la \VUgras{kupran} \textsuperscript{(44)} \VUgras{fadenon}
de la telefono.

%%K t05-tit-7 sub
\VUpk{10.2pt}{40}{La ĝardenkultivado.}
\VUsaut{3.0mm}

%%K t05-01-43 p
\textsc{La onklo}. --- En la fundo de la \VUgras{korto} \textsuperscript{(45)},
apud la \VUgras{krado} \textsuperscript{(47)} kaj la \VUgras{pordego}
\textsuperscript{(48)} vi vidas la \VUgras{ĝardeniston}
\textsuperscript{(129)}. Li preparas la \VUgras{ĝardeneton}
\textsuperscript{(49)}. Li fosis la teron per sia \VUgras{fosilo}
\textsuperscript{(131)}, li semas semojn kaj rastas per la \VUgras{rastilo}
\textsuperscript{(130)}. Poste, per sia \VUgras{verŝilo} \textsuperscript{(132)},
li akvumos la \VUgras{bedojn} \textsuperscript{(150)} kaj la plantoj kreskados.

%%K t05-01-44 p
La \VUgras{ĉevalservisto} \textsuperscript{(133)} kiu ĵus flegis la ĉevalon kaj
donis al ĝi nutraĵon purigas la \VUgras{veturilon} \textsuperscript{(52)} per
spongo, \VUgras{manpumpilo} \textsuperscript{(134)} kaj sitelo da akvo. Poste
li reenirigos ĝin en la veturilejon kaj fermos la pordon per la dika
\VUgras{riglilo} \textsuperscript{(53)}.

%%K t05-01-45 p
Jen vi kontenta, mi opinias, vi havis sufiĉajn klarigojn.

%%K t05-01-46 p
\textsc{Johano}. --- Jes, onklo mia, sed volonte mi vidus la internon de la
domo.

%%K t05-01-47 p
\textsc{La onklo}. --- Tio okazos morgaŭ. Sinjoro Rufo klarigos al vi la planon
kaj montros al vi la nomon de ĉiu ĉambro.

%%K t05-tit-8 sub
\VUsaut{3.0mm}
\VUpk{10.2pt}{40}{Interna aranĝo de la domo.}
\VUsaut{2.4mm}

%%K t05-apar-2 apar
{\centering\textit{(Vidu la planon.)}\par}
\VUsaut{2.4mm}

%%K t05-01-48 p
\textsc{La arkitekto}. --- Venu, Johano, mi tuj igos vin esplori la diversajn
partojn de la domo.

%%K t05-01-49 p
Ni eniru la \VUgras{teretaĝon} \textsuperscript{(7)}. Jen la butono de la
\VUgras{sonorileto} \textsuperscript{(11)}. Ni supreniras la tri
\VUgras{ŝtupojn} \textsuperscript{(14)} de la \VUgras{perono}
\textsuperscript{(9)}, ni transiras la \VUgras{sojlon} \textsuperscript{(10)}
de la enir-\VUgras{pordo} \textsuperscript{(8)} kaj jen ke ni estas ĉe la piedo
de la \VUgras{ŝtuparo} \textsuperscript{(13)}.

%%K t05-01-50 p
\textsc{Johano}. --- Tio estas la koridoro.

%%K t05-01-51 p
\textsc{La arkitekto}. --- Ne, tio estas la \VUgras{vestiblo}
\textsuperscript{(12)}. La \VUgras{koridoro} \textsuperscript{(a)} estas ĉe la
ekstremo. Dekstre, jen la \VUgras{salono} \textsuperscript{(b)} kaj la
\VUgras{manĝoĉambro} \textsuperscript{(c)}, kontraŭ la manĝoĉambro estas la
\VUgras{kuirejo} \textsuperscript{(d)} kaj la \VUgras{servistejo}
\textsuperscript{(e)}, poste antaŭe la \VUgras{skribĉambro}
\textsuperscript{(f)}. La \VUgras{verando} \textsuperscript{(23)} kun sia
\VUgras{vitrita pordo} \textsuperscript{(25)} estas apud la salono.

%%K t05-01-52 p
Ni tuj supreniros la ŝtuparon. Subtenu vin per la \VUgras{balustrado}
\textsuperscript{(15)} kaj kalkulu la ŝtupojn. Ili estas dek kvar. Jen la
\VUgras{placeto} \textsuperscript{(m)}, ni estas sur la unua \VUgras{etaĝo}
\textsuperscript{(27)}.

%%K t05-tit-9 sub
\VUsaut{3.0mm}
\VUpk{10.2pt}{40}{La unua etaĝo.}
\VUsaut{3.0mm}

%%K t05-01-53 p
Kontraŭ la ŝtuparo, estas la \VUgras{necesejo} \textsuperscript{(20)} kaj la
\VUgras{tualetejo} \textsuperscript{(j)}. Dekstre bela \VUgras{dormoĉambro}
\textsuperscript{(g)} kun du fenestroj ĉe la \VUgras{fasado}
\textsuperscript{(1)} de la domo. Maldekstre, estas la \VUgras{banejo}
\textsuperscript{(1)} kaj la \VUgras{ĉambro por amikoj} \textsuperscript{(i)}
kun granda \VUgras{alkovo} \textsuperscript{(h)}: Apud la dormoĉambro estas la
\VUgras{ĉambro por la infanoj} \textsuperscript{(k)}. Ĝi situas ĉe vasta
\VUgras{teraso} \textsuperscript{(21)}. Sub ĉi tiu teraso estas alia necesejo
(20) kaj, ĉe la muro, jen la \VUgras{sonorilo} \textsuperscript{(22)} kiu
sonas por la vespermanĝo.

%%K t05-01-54 p
\textsc{Johano}. --- Ni iru sur la \VUgras{duan etaĝon}.

%%K t05-01-55 p
\textsc{La arkitekto}. --- Ne ekzistas dua etaĝo. Sub la tegmento estas nur
\VUgras{mansardoj} \textsuperscript{(33)} kaj la grenejo (34). Ni finis la
esploron, ni povas malsupreniri.

%%K t05-01-56 p
\textsc{Johano}. --- Dankon, sinjoro, tio estas tre interesa. Mi tuj rakontos
al mia patro ĉion kion mi vidis.
\VUsaut{4.0mm}
\VUfilet{22mm}
